\section{Lean definitions and a path to complete object formalization}
\label{sec:formalization}

The formalization has three different products. The distinction prevents
the common error of replacing an object-theory proof by a theorem of
the surrounding prover.

\begin{enumerate}
\item \texttt{AffineBounds.lean} contains eighteen proved theorems in
Lean about natural-number recurrences and finite sums. It is compiled.
\item The blueprint modules contain compiled definitions of rational
data, algorithms, guards and named target propositions. Defining a
proposition does not prove it. These modules introduce no new axiom
asserting their targets.
\item \texttt{ObjectAdapter.lean} is an uncompiled integration outline
for actual \texttt{Prf PRAAx} derivations, at the pinned ProveIt
revision. Its proof fields must be populated by ordinary iPRA proof
objects.
\end{enumerate}

The arithmetic proofs can be checked independently of Mathlib and
ProveIt. The final object export requires the existing repository
arithmetic library and the additional instances described here.

\subsection{Finite rational definitions}

The coupled blueprint uses the host definitions
\begin{lstlisting}
abbrev Vec (d : Nat) := Fin d -> Rat
abbrev Mat (d : Nat) := Fin d -> Fin d -> Rat

structure Data (d : Nat) where
  matrix : Mat d
  weights : Vec d
  reciprocalGap : Nat

def contraction (c : Data d) : Rat :=
  1 - 1 / (c.reciprocalGap : Rat)
\end{lstlisting}

It defines finite sums and maxima by folds, and defines
\texttt{weightedNorm}, \texttt{InBox}, \texttt{SmallGain},
\texttt{Diameter} and \texttt{ActionContract} exactly as in the
human theorem. The precision and iteration functions are literal
natural-number expressions, and \texttt{orbit} is a primitive
recursion. Other definitions include vector residuals, finite Neumann
prefixes, automatic weights, finite state histories, delay guards,
and update-coverage guards.

The named propositions include
\nolinkurl{WeightedMaximumTarget}, \nolinkurl{ResidualSeparationTarget},
\nolinkurl{ProducerTarget}, \nolinkurl{NeumannCertificateTarget},
\nolinkurl{AutomaticWeightsTarget}, \nolinkurl{PosteriorErrorTarget},
\nolinkurl{DelayedRunTarget} and \nolinkurl{AsynchronousProducerTarget}.
They expose all assumptions. For example, the producer target is
the conjunction of box membership, the finite residual estimate,
and regularity of the explicitly defined sampler. It does not
conclude that some unspecified function exists.

These host functions are specifications for the coded backend, not
an assertion that every Lean function is PR. The object backend
replaces \texttt{Fin d -> Rat} by a length-\(d\) sequence code,
and each evaluator parameter by a registered PR body with its
representation equations. The definitions in \cref{sec:arithmetic}
specify the full numerical representation needed for that step.

\clearpage
\subsection{Core theorem dependency order}

\begin{longtable}{@{}>{\raggedright\arraybackslash}p{0.21\textwidth}>{\raggedright\arraybackslash}p{0.37\textwidth}>{\raggedright\arraybackslash}p{0.34\textwidth}@{}}
\toprule
Target & Definitions and prerequisite lemmas & Proof to implement\\
\midrule
\endhead
Weighted comparison & Rational order; finite maximum with attaining index;
nonnegative matrix multiplication & The positive-maximum contradiction in
\cref{lem:weighted-maximum}, after denominator clearing\\
Residual separation & Weighted triangle inequality; mixed action contract &
Three intermediate vectors and cancellation of \((s-1)d\)\\
Geometric clock & Natural powers; Bernoulli bound; rational division &
QF induction on \(q^n(1+n/s)\le1\), then on block count\\
Regular producer & PR orbit; box bit; diameter; affine envelope &
QF orbit invariant followed by substitution of \(P(k),N(k)\)\\
Neumann enclosure & Matrix powers; finite sum identity &
QF induction on the prefix identity, then weighted comparison\\
Inverse sampler & Induced matrix norm; submultiplicativity; powers &
Finite two-sided telescoping and geometric tail\\
Polynomial chart & Sparse polynomial evaluation, substitution and
formal derivative & Monomial factorization, exact coefficient cancellations,
then the explicit radius inequalities\\
Delayed execution & Finite history; read-time table; update mask &
QF history-window invariant and the bounded coverage bit\\
Refinement coherence & Coded restrictions; nonexpansiveness; residuals &
One rational cancellation, then finite telescoping over model levels\\
Lattice sampler & Dyadic floor; uniform polynomial evaluator; prefix &
Box, contraction, commutation and rounding lemmas, followed by substitution\\
\bottomrule
\end{longtable}

\subsection{The actual induction formulas}

Fix all input parameters and one finite horizon. The following are
the essential invariants, written with rational abbreviations:
\begin{align*}
 \mathcal I_{\mathrm{orbit}}(t):\quad&
 x_t\in D\ \wedge\ x_{t+1}\in D\ \wedge\
 d_w(x_t,x_{t+1})\le q^tR+2s2^{-p};\\
 \mathcal I_{\mathrm{matrix}}(t):\quad&
 A^tw\le q^tw\ \wedge\
 A\sum_{j<t}A^jr+r=\sum_{j<t}A^jr+A^tr;\\
 \mathcal I_{\mathrm{round}}(t):\quad&
 v_t,z_t\in D_n\ \wedge\ d_n(v_t,z_t)\le2h;\\
 \mathcal I_{\mathrm{tower}}(t):\quad&
 d_n(y_n,P_{n,n+t}y_{n+t})\le
       \sum_{j=n}^{n+t-1}e_j.
\end{align*}
Every coordinate conjunction is bounded by a dimension, and every
sum is a PR term. Replace the bounded predicates by their
characteristic-bit equations and clear positive denominators.
The resulting formulas are quantifier-free. If hypotheses are
supplied only through a finite trace, conjoin the corresponding
prefix-validation bit or use its implication form (6).

For a delayed block, the state includes the finite history, the
block counter, and a Boolean mask for coordinates already updated
in that block. The invariant says that all entries in the last
delay window satisfy the old bound, that updated coordinates
satisfy the improved bound, and that all states remain in the
Cartesian box. All time quantifiers range over the finite stored
history. The coverage hypothesis is another bounded bit.

For the tower, compute \(P_{n,m}(z)\) by a loop of length
\(m-n\) that applies \(P_{m-1},P_{m-2},\ldots,P_n\) in
that order. Its level counter decreases by truncated subtraction,
and its vector code is the current state. This is an ordinary
first-order PR implementation of restriction composition; no
higher-type recursion on arbitrary maps is required.

\subsection{Clearing rational denominators correctly}

The compiled integer envelope theorem is not applied to rational
errors without an encoding step. There are two convenient choices.
One can carry numerator and denominator fields through the proof
and use the rational comparison bit directly. Alternatively, for
one finite horizon \(T\), compute a common positive multiple of
all denominators appearing in the data and the finite trace. A
product is sufficient; a least common multiple is unnecessary.
More precisely, let \(e_t\ge0\) be the rational error for
\(t\le T\), and let \(f_t\ge0\) be the rational forcing term
for \(t<T\). Let \(L_T>0\) be that denominator product, including
the denominator of every \(e_t\) through \(T\) and every
\(f_t\) before \(T\). Then
\[
 a_t=L_Te_t\quad(t\le T),\qquad
 r_t=L_Tf_t\quad(t<T)
\]
are nonnegative integers. In the coded implementation they are PR
integer terms, obtained by integer multiplication and division;
divisibility by each original denominator follows from the product
definition. The terms retain \(T\) and the original input as
numerical parameters. Multiplication by the positive \(L_T\)
preserves each recurrence inequality on this finite horizon.

The guarded-loop interface needs these inequalities only for
\(t<T\). To instantiate the compiled host theorem in its
all-natural-index form, one may instead extend \(a_t\) by zero
for \(t>T\), and \(r_t\) by zero for \(t\ge T\).
The transition at \(T\) then has zero left side, and every later
transition is \(0\le0\). This explicit totalization preserves
the value at \(T\) and requires no common denominator for an
infinite sequence.

For the scalar contraction route, the coefficient is exactly
\((s-1)/s\), so the cleared recurrence has the form
\(s a_{t+1}\le(s-1)a_t+r_t\). This matches the compiled clock
and the constant-forcing lemmas. If rational matrix entries are
cleared directly, the resulting numerator and denominator need
not differ by one. Use the general
\texttt{coupled\_envelope} with arbitrary \(c,d\), or first
pass to the rational weighted maximum and then use the scalar
route. This distinction avoids imposing an accidental integrality
restriction on the matrix theorem.

The denominator product may be large, but is a PR function of the
finite trace. The proof does not require a complexity bound from
a stronger arithmetic theory. In the dyadic lattice implementation,
the working denominators are particularly simple because each
state is rounded back to its declared dyadic grid.

\subsection{Object-level packaging and discharged assumptions}

The adapter's \texttt{Package} contains actual PR terms for initial
state, step and fuel; a bounded invariant syntax; a guard and
result formula; and a \texttt{GuardedLoop.Loop} certificate.
The certificate requires actual object proofs of the base,
preservation, and final implication. The export has the form
\begin{lstlisting}
Gamma |- [PRAAx] guard(u)
------------------------------------------
Gamma |- [PRAAx] result(Run(init, step, fuel(u), u), u)
\end{lstlisting}
The displayed notation is explanatory; the file uses the repository's
literal judgment and constructors.

For the general approximate-action theorem, the box and mixed-action
contracts remain assumptions in the sequent. For the polynomial
frontend they are discharged by coefficient algebra. For the
infinite lattice, the concrete polynomial, weight, floor and
boundary estimates discharge every such analytic assumption.
The remaining premises in an object proof are then only its
ordinary numerical input conditions, such as positive dimension.
The construction should export the actual universally quantified
rational-code sentences listed in \cref{sec:lattice}.

No global principle ``every accepted proof trace is true'' is
postulated. The existing \texttt{Prf.trans}, substitution and
loop theorems compose proof values supplied at the relevant
formulas. The elementary arithmetic obligations still need to
be encoded and checked in that calculus.

\subsection{Semantic interpretation is an additional theorem}

The principal iPRA outputs are finite rational estimates. In an
ordinary semantic layer one may interpret their regular samplers
as real vectors, coordinate functions or inverse-limit objects.
The norm estimates and residuals prove that these interpreted
objects satisfy the advertised equations. That interpretation
does not strengthen the object induction rule, nor does a proof
about external real numbers automatically return an object
derivation.

A complete Lean development should therefore expose two named
endpoints for each application: its literal coded iPRA sentence
with a \texttt{Prf PRAAx} proof, and its semantic adequacy theorem
in the chosen represented-object library. The present packet
contains the human proofs and the detailed construction path;
it does not claim those two final layers have already been
implemented in full.
